\section{Best Practices} \label{appendix:best}

\subsection{General Principle}

This paper introduces two algorithms: \ours and \varn[n], with a key focus on selecting the appropriate algorithm based on task conditions.

Empirical evidence from the open-source community suggests that \varn[n] is particularly effective for sample filtering or more complex scenarios. For instance, in the multi-turn tool-calling setting, a high proportion of void (i.e., non-informative or incorrect) samples can destabilize training. In such cases, incorporating a baseline by subtracting the intra-group mean reward significantly improves training stability by effectively filtering out void samples. Moreover, this automatic reward reshaping reduces the need for designing complex reward structures. For example, \varn[n] supports both 0/1 and -1/1 reward schemes, while \ours performs best with symmetric rewards, such as -1/1 in RLVR tasks.

In contrast, for general-domain tasks where prompt diversity and efficiency are paramount, or for tasks where obtaining multiple reward signals for distinct responses is challenging—such as training with Process-Supervised Reward Models~(PRMs) or online real-time sampling—we recommend using \textbf{\ours ($k=1$)}. Our experiments (Section 4.1) show that it provides superior OOD generalization in these settings.

\subsection{Third-Party Validation}
The core principle of \ours, global advantage normalization, has been independently validated and adopted in several \textbf{subsequent} large-scale reasoning systems, confirming its stability and effectiveness.

\paragraph{LitePPO}~\citep{liu2025part}, which effectively combines \varn with a token-level loss, conducted experiments demonstrating that the global standard deviation is superior to the local standard deviation used in GRPO. Their results indicate more stable training and better generalization, which aligns with our findings.

\paragraph{ScaleRL}~\citep{khatri2025art} performed a detailed ablation study on advantage estimation methods in large-scale (16,000 GPU-hour) experiments. They directly compared batch-level normalization in \ours with prompt-level normalization in GRPO. Their findings concluded that batch-level normalization was "slightly superior in both compute efficiency and final performance," confirming the advantages of global normalization at scale.

\paragraph{DLER}~\citep{liu2025dler} found that when using truncation to control the output length of an LLM, batch-wise (global) normalization remains stable while group-wise (local) normalization shows declining accuracy. The study provides further evidence that global normalization is more robust to variations in training conditions and reward landscapes.